\subsection{Baselines and what the map predicts}
Table~\ref{tab:base} reports reconstruction RMSE (mV) of the unobserved target leads under a
\emph{single, identical protocol} for all methods: per-timepoint waveform RMSE on the same
$\NTestFair$ held-out fold-10 records, same segment indices, record-bootstrap CIs. This is
the fair comparison---an earlier draft scored the classical baselines on segment-mean vectors
and the neural baseline on waveforms, which is not comparable. We include a \emph{strong
neural baseline}: an arbitrary-mask conditional 1-D U-Net trained with random lead masks (so
one model serves any configuration, no target-lead leakage), the masked-reconstruction
architecture the reduced-lead literature uses \cite{gradowski2025reveals}.

On a dipole-spanning set the methods form a monotone ladder---prior-mean $\to$ dipolar
(Tier~I) $\to$ per-segment ridge (Tier~I+II) $\to$ U-Net---confirming there is recoverable
non-dipolar content (the calibrated Tier~II layer), yet the strong neural model's margin over
a simple ridge is small (QRS $\FairSpanUnetQRS$ vs.\ $\FairSpanRidgeQRS$\,mV on
$\{$I,II,V1,V3,V5$\}$). The decisive result is on \textbf{limb-6}: \emph{every} reconstructor
degrades sharply, and the U-Net---despite its capacity---reaches only QRS $\FairLimbUnetQRS$,
$\FairUnetLimbSpanRatio\times$ its spanning-set error, still far above what it achieves when
the target is identifiable. No amount of model capacity crosses the identifiability ceiling
the map's $\eta{>}0$ warning marks out \emph{a priori}: limb$\to$precordial morphology is not
recoverable, so a low reconstruction error there would be spurious.

\begin{table}[t]
\centering\small
\caption{Reconstruction RMSE (mV) of target leads, QRS segment, under one identical
per-timepoint protocol ($\NTestFair$ test records). ``dipolar'' = Tier~I; ``ridge'' =
per-segment ridge (Tier~I+II); ``U-Net'' = strong arbitrary-mask neural baseline. Best in
bold. On limb-6 even the neural baseline cannot recover precordial morphology, as $\eta{>}0$
predicts. Numbers auto-generated from \texttt{results/fair\_baselines.json}.}
\label{tab:base}
% AUTO-GENERATED by paper/emit_baseline_table.py -- do not edit by hand.
\begin{tabular}{lcccc}
\toprule
config & prior-mean & dipolar & ridge & U-Net \\
\midrule
$\{$I,II,V2$\}$ & 0.471 & 0.234 & 0.220 & \textbf{0.181} \\
$\{$I,II,V1,V3,V5$\}$ & 0.484 & 0.186 & 0.168 & \textbf{0.164} \\
limb-6 & 0.484 & 0.362 & 0.346 & \textbf{0.239} \\
\bottomrule
\end{tabular}

\end{table}

\textbf{Empirical, not physical, dipole.} On real PTB-XL the estimated $\bm M_s$ shares two
directions with the classical inverse-Dower vectorcardiographic space (principal angles
$2$--$6^\circ$) but its third direction differs materially (max angle $43$--$55^\circ$,
bootstrap CIs excluding small angles). We therefore treat $\bm M_s$ as an \emph{empirical
rank-3 spatial subspace}, not a physical cardiac dipole; the earlier synthetic
``matches inverse-Dower'' check used inverse-Dower to generate the data and is not evidence
on real ECG. All certificate quantities are stated with respect to this estimated subspace.
